% Options for packages loaded elsewhere
\PassOptionsToPackage{unicode}{hyperref}
\PassOptionsToPackage{hyphens}{url}
\documentclass[
  11pt,
]{article}
\usepackage{xcolor}
\usepackage[margin=25mm]{geometry}
\usepackage{amsmath,amssymb}
\setcounter{secnumdepth}{-\maxdimen} % remove section numbering
\usepackage{iftex}
\ifPDFTeX
  \usepackage[T1]{fontenc}
  \usepackage[utf8]{inputenc}
  \usepackage{textcomp} % provide euro and other symbols
\else % if luatex or xetex
  \usepackage{unicode-math} % this also loads fontspec
  \defaultfontfeatures{Scale=MatchLowercase}
  \defaultfontfeatures[\rmfamily]{Ligatures=TeX,Scale=1}
\fi
\usepackage{lmodern}
\ifPDFTeX\else
  % xetex/luatex font selection
\fi
% Use upquote if available, for straight quotes in verbatim environments
\IfFileExists{upquote.sty}{\usepackage{upquote}}{}
\IfFileExists{microtype.sty}{% use microtype if available
  \usepackage[]{microtype}
  \UseMicrotypeSet[protrusion]{basicmath} % disable protrusion for tt fonts
}{}
\makeatletter
\@ifundefined{KOMAClassName}{% if non-KOMA class
  \IfFileExists{parskip.sty}{%
    \usepackage{parskip}
  }{% else
    \setlength{\parindent}{0pt}
    \setlength{\parskip}{6pt plus 2pt minus 1pt}}
}{% if KOMA class
  \KOMAoptions{parskip=half}}
\makeatother
% definitions for citeproc citations
\NewDocumentCommand\citeproctext{}{}
\NewDocumentCommand\citeproc{mm}{%
  \begingroup\def\citeproctext{#2}\cite{#1}\endgroup}
\makeatletter
 % allow citations to break across lines
 \let\@cite@ofmt\@firstofone
 % avoid brackets around text for \cite:
 \def\@biblabel#1{}
 \def\@cite#1#2{{#1\if@tempswa , #2\fi}}
\makeatother
\newlength{\cslhangindent}
\setlength{\cslhangindent}{1.5em}
\newlength{\csllabelwidth}
\setlength{\csllabelwidth}{3em}
\newenvironment{CSLReferences}[2] % #1 hanging-indent, #2 entry-spacing
 {\begin{list}{}{%
  \setlength{\itemindent}{0pt}
  \setlength{\leftmargin}{0pt}
  \setlength{\parsep}{0pt}
  % turn on hanging indent if param 1 is 1
  \ifodd #1
   \setlength{\leftmargin}{\cslhangindent}
   \setlength{\itemindent}{-1\cslhangindent}
  \fi
  % set entry spacing
  \setlength{\itemsep}{#2\baselineskip}}}
 {\end{list}}
\usepackage{calc}
\newcommand{\CSLBlock}[1]{\hfill\break\parbox[t]{\linewidth}{\strut\ignorespaces#1\strut}}
\newcommand{\CSLLeftMargin}[1]{\parbox[t]{\csllabelwidth}{\strut#1\strut}}
\newcommand{\CSLRightInline}[1]{\parbox[t]{\linewidth - \csllabelwidth}{\strut#1\strut}}
\newcommand{\CSLIndent}[1]{\hspace{\cslhangindent}#1}
\setlength{\emergencystretch}{3em} % prevent overfull lines
\providecommand{\tightlist}{%
  \setlength{\itemsep}{0pt}\setlength{\parskip}{0pt}}
\usepackage{bookmark}
\IfFileExists{xurl.sty}{\usepackage{xurl}}{} % add URL line breaks if available
\urlstyle{same}
\hypersetup{
  hidelinks,
  pdfcreator={LaTeX via pandoc}}

\author{}
\date{}

\begin{document}

\section{Fixed-horizon information bounds for deterministic dynamics in
OU measurement
error}\label{fixed-horizon-information-bounds-for-deterministic-dynamics-in-ou-measurement-error}

The results below specialize classical Gaussian shift theory to the
observation model. Equivalence-based impossibility is established
methodology, including Zhang (2004, Section 3); the contribution here is
an explicit dynamical sensitivity norm and its link to weak residuals
and acquisition choices.

Continuous-record precision and efficient discrete designs also have
precedents in correlated regression {[}1{]}. Local ODE sensitivities
serve as regressors in that setting.

\subsection{Observation experiment and Cameron--Martin
norm}\label{observation-experiment-and-cameronmartin-norm}

Observe \(Y(t)=\mu_\theta(t)+\varepsilon(t)\) on \([0,T]\), where the
deterministic mean is in \(H^1([0,T];\mathbb R^d)\). Every parameter
uses the same known stationary OU covariance. Coordinates have positive
variances \(\sigma_j^2\) and correlation times \(\tau_j>0\), and are
independent. The RKHS is \(H^1\) as a set, with norm

\[
\|g\|_{\mathcal H}^2=\sum_{j=1}^d\left[\frac{g_j(0)^2}{\sigma_j^2}+\frac{\tau_j}{2\sigma_j^2}\int_0^T\left(g'_j+\frac{g_j}{\tau_j}\right)^2dt\right].
\]

To verify this directly, write
\(d\varepsilon_j=-\varepsilon_jdt/\tau_j+\sqrt{2\sigma_j^2/\tau_j}\,dW_j\),
with an independent stationary Gaussian initial value. The invertible
path transformation
\(z\mapsto(z(0),z(t)-z(0)+\tau_j^{-1}\int_0^tz(s)ds)\) maps the
experiment to its initial normal variable and scaled Brownian path.
Translation by \(g\) costs the initial normal shift plus Brownian drift
\(g'+g/\tau_j\). The Cameron--Martin theorem gives the stated norm.
Expanding its square gives the symmetric form

\[
\|g\|_{\mathcal H}^2=\sum_j\left[\frac{\tau_j}{2\sigma_j^2}\int_0^T(g'_j)^2dt+\frac1{2\sigma_j^2\tau_j}\int_0^Tg_j^2dt+\frac{g_j(0)^2+g_j(T)^2}{2\sigma_j^2}\right].
\]

The endpoint terms are essential for stationary initial noise. Assuming
zero initial noise would define a different experiment.

\subsection{Theorem: impossibility for an entire parameter
family}\label{theorem-impossibility-for-an-entire-parameter-family}

For any two distinct parameters \(\theta,\eta\) with means in \(H^1\),
their continuous-path laws are equivalent. Let
\(\delta=\|\theta-\eta\|\) and
\(d_{\theta\eta}=\|\mu_\theta-\mu_\eta\|_{\mathcal H}\). Every estimator
\(\widehat\theta_h\) based on any finite collection of evaluations of
the same path satisfies

\[
\max_{v\in\{\theta,\eta\}}P_v\left(\|\widehat\theta_h-v\|\geq\delta/2\right)\geq\Phi(-d_{\theta\eta}/2)>0.
\]

Consequently no sequence of denser-grid estimators is consistent in
probability at every parameter of a nontrivial family. This assertion
permits an estimator that is consistent at an isolated parameter, for
example a constant estimator.

\textbf{Proof.} A shift between the means belongs to the common RKHS, so
the Cameron--Martin translation theorem gives mutual absolute
continuity. Under the first law the log likelihood ratio of the second
law is normal with mean \(-d_{\theta\eta}^2/2\) and variance
\(d_{\theta\eta}^2\). Thus the total variation distance of the path laws
is \(2\Phi(d_{\theta\eta}/2)-1\). Evaluation of a path cannot increase
total variation. Any estimator defines a nearest-parameter test. A test
error can occur only when the estimator's error relative to the
generating parameter is at least \(\delta/2\). The sum of errors of
every test is at least one minus total variation, or
\(2\Phi(-d_{\theta\eta}/2)\). At least one of the two estimation error
probabilities is at least half this sum. The lower bound is independent
of the grid, which excludes simultaneous consistency. The same reasoning
and \(E[Z^2]\geq r^2P(|Z|\geq r)\) yield

\[
\max_{v\in\{\theta,\eta\}}E_v\|\widehat\theta_h-v\|^2\geq\frac{\delta^2}{4}\Phi(-d_{\theta\eta}/2).
\]

This risk statement allows infinite moments and needs neither
unbiasedness nor a particular weak estimator. It also covers randomized
estimators by adjoining independent randomization to both experiments.
For \(d_{\theta\eta}=0\), both means define the same law and the bound
remains valid. \(\square\)

\subsection{Fisher information and weak
efficiency}\label{fisher-information-and-weak-efficiency}

Suppose \(\theta\mapsto\mu_\theta\) is differentiable in the RKHS norm.
For sensitivities \(s_r=\partial_{\theta_r}\mu_\theta\), the
continuous-record Fisher information is

\[
\mathcal I_{rs}(\theta)=\langle s_r,s_s\rangle_{\mathcal H}.
\]

For a regular unbiased estimator of a scalar parameter,
\(\operatorname{Var}(\widehat\theta)\geq\mathcal I(\theta)^{-1}\). With
nuisance parameters the appropriate bound uses the corresponding element
of the inverse information matrix, when nonsingular. Finite information
does not imply that a finite-horizon maximum likelihood estimator
attains this bound.

For \(\mu_a(t)=x_0e^{at}\) with known \(x_0\), sensitivity
\(s(t)=x_0te^{at}\) and \(s(0)=0\), one obtains

\[
\mathcal I_T(a)=\frac{\tau x_0^2}{2\sigma^2}\int_0^T e^{2at}\left[1+(a+\tau^{-1})t\right]^2dt.
\]

At a uniform grid, \(\rho=e^{-h/\tau}\), the exact information is

\[
\mathcal I_h(a)=\frac{s(0)^2}{\sigma^2}+\frac{\sum_{n=1}^N[s(t_n)-\rho s(t_{n-1})]^2}{\sigma^2(1-\rho^2)}\ \longrightarrow\ \mathcal I_T(a).
\]

For nested grids it increases monotonically by data processing. The
convergence also follows directly from
\(s(t_n)-\rho s(t_{n-1})=h(s'+s/\tau)+o(h)\) and
\(1-\rho^2=2h/\tau+o(h)\).

For the scalar weak ratio, define \(D_\phi=\int\phi\mu_a\ne0\). Its
first-order noise influence is \(\ell(t)=(-\phi'(t)-a\phi(t))/D_\phi\).
Differentiating the weak identity gives \(\int\ell s=1\). Let
\(V_\phi=\iint\ell(u)C(u,v)\ell(v)dudv\). The RKHS representer of this
linear functional is \(C\ell\), so Cauchy--Schwarz gives

\[
1=\langle C\ell,s\rangle_{\mathcal H}^2\leq V_\phi\mathcal I_T(a),\qquad \operatorname{Eff}_\phi=\frac1{\mathcal I_T(a)V_\phi}\leq1.
\]

Here the squared inner product equals one; the preceding display uses
that normalization. \(V_\phi\) is the variance of a Gaussian
linearization. The unguarded ratio itself can have infinite second
moment, so its empirical squared errors are not estimators of \(V_\phi\)
and must not be described as an attained Cramér--Rao variance.

\subsection{Initial-state annihilation and a stricter weak influence
bound}\label{initial-state-annihilation-and-a-stricter-weak-influence-bound}

Let \(S=\partial_\theta x\) and \(Z=\partial_u x\) for active
vector-field coefficients and initial state \(u\). Their variational
equations are \(S'=JS+F\) and \(Z'=JZ\), where \(J=\partial_xf\) and
\(F=\partial_\theta f\). For compactly supported test functions, the
first-order joint residual influence is
\(A_k(t)=-\phi'_k(t)I-\phi_k(t)J(t)\). Integration by parts yields
\(\int A_kZ=0\) and \(\int A_kS=\int\phi_kF\). With a full-rank clean
weak design and a left inverse \(P_0\), the normalized coefficient
influence \(\ell=P_0A\) therefore satisfies

\[
\int\ell Z=0,\qquad\int\ell S=I.
\]

Write the Gaussian sensitivity score vectors as \(U_\theta,U_u\). They
have joint information matrix given by the RKHS Gram matrix of \(S,Z\).
When its blocks and Schur complement are positive definite, the
efficient score is

\[
U_{\mathrm{eff}}=U_\theta-\mathcal I_{\theta u}\mathcal I_{uu}^{-1}U_u,
\qquad
\mathcal I_{\mathrm{eff}}=\mathcal I_{\theta\theta}-\mathcal I_{\theta u}\mathcal I_{uu}^{-1}\mathcal I_{u\theta}.
\]

For \(L_\varepsilon=\int\ell\varepsilon\), the score covariance identity
gives \(\operatorname{Cov}(L_\varepsilon,U_{\mathrm{eff}})=I\). Thus

\[
\operatorname{Cov}\big(L_\varepsilon-\mathcal I_{\mathrm{eff}}^{-1}U_{\mathrm{eff}}\big)
=\operatorname{Cov}(L_\varepsilon)-\mathcal I_{\mathrm{eff}}^{-1}\succeq0.
\]

The Schur-complement covariance bound is classical Gaussian regression
geometry {[}1{]}. The weak restriction is its mandatory annihilation of
initial-state sensitivities, even if the actual initial state is
supplied. This bound concerns the Gaussian influence, not the unguarded
nonlinear estimator's moments. It is not an attainment theorem for a
chosen finite window family.

For \(x=x_0e^{at}\), initial-state sensitivity is \(q=e^{at}\), while
coefficient sensitivity is \(s=x_0te^{at}\). The projected information
is
\(\mathcal I_{aa}-\langle s,q\rangle_{\mathcal H}^2/\|q\|_{\mathcal H}^2\).
Independent deterministic numerical verification is recorded in the
geometry execution manifest.

\subsection{Acquisition regimes and heterogeneous
recordings}\label{acquisition-regimes-and-heterogeneous-recordings}

For a constant mean \(\mu_\theta(t)=\theta\), the exact continuous
information is \(\sigma^{-2}[1+T/(2\tau)]\). The quantity
\(1+T/(2\tau)\) is its information-equivalent number of independent
observations of variance \(\sigma^2\). This identifies the often-used
\(T/\tau\) scale, including its factor of two, in a specific model.

For a decaying exponential \(a<0\), \(\mathcal I_T(a)\) converges to a
finite limit as \(T\to\infty\); a longer near-zero tail does not
maintain parameter excitation. At \(a=0\), information instead grows as
\(\tau x_0^2[ T+T^2/\tau+T^3/(3\tau^2)]/(2\sigma^2)\). Therefore a
universal duration rule based only on \(T/\tau\) is inappropriate for
dynamical coefficients.

Independent recordings add information: identical records yield
\(M\mathcal I_T\), while records with known differing initial conditions
yield \(\sum_{m=1}^M\mathcal I_{T,m}\). An acquisition decision must
therefore compare model sensitivities, physical covariance, and cost. A
large \(\mathcal I_h/\mathcal I_T\) diagnoses how much information a
grid captures in the known Gaussian model; it is neither a
support-recovery guarantee nor an operational estimate when
sensitivities are unknown.

For weak estimation with different initial conditions and a shared
polynomial coefficient vector, every clean pair satisfies
\(b_m^0=G_m^0\xi\). Corrected weak matrices have conditional
expectations \(E[G_m\mid x_m]=G_m^0\) and \(E[b_m\mid x_m]=b_m^0\) up to
deterministic quadrature error. Under independent identically
distributed recordings, integrability and full rank of \(E[G_m^0]\) on
the supplied support, the law of large numbers implies consistency of
least squares applied to \(\overline G_M,\overline b_M\), first for the
population finite-grid identity and subsequently as quadrature error
vanishes. The full-rank condition is substantive: opposing excitation
can cancel when matrices are averaged. Stacking designs instead involves
\(E[G_m^TG_m]\) and \(E[G_m^Tb_m]\), which generally retain
errors-in-variables terms. Averaging sensor paths before evaluating
nonlinear features is invalid for differing initial conditions because
\(f(E[x_m])\) need not equal \(E[f(x_m)]\).

\subsection{Scope}\label{scope}

All equivalence and information conclusions use a fixed finite horizon
and common positive OU covariance across parameters. They need not hold
for noiseless coordinates, parameter-dependent singular covariances,
repeated independent recordings with \(M\to\infty\), or
increasing-domain experiments with persistent excitation. The classical
measure-equivalence mechanism is attributed to its predecessors; its
specialization should not be advertised as a new Gaussian-measure
theorem.

\protect\phantomsection\label{refs}
\begin{CSLReferences}{0}{0}
\bibitem[\citeproctext]{ref-dette2016}
\CSLLeftMargin{{[}1{]} }%
\CSLRightInline{Dette H, Pepelyshev A, Zhigljavsky A. Optimal designs in
regression with correlated errors. \emph{The Annals of Statistics} 44
(2016) 113--152. \url{https://doi.org/10.1214/15-AOS1361}.}

\end{CSLReferences}

\end{document}
